\section*{Chapter \ref{ch:g11:greedy-algorithms} --- Greedy Algorithms}

\begin{solution}{exo:g11:greedy-algorithms:1}
$50$ leaves $37$; $20$ leaves $17$; $10$ leaves $7$; $5$ leaves $2$; $2$ leaves $0$: five coins, $50 + 20 + 10 + 5 + 2$.
\end{solution}

\begin{solution}{exo:g11:greedy-algorithms:2}
Greedy gives $4 + 1 + 1$, three coins; $3 + 3$ needs two.
\end{solution}

\begin{solution}{exo:g11:greedy-algorithms:3}
By end hour: P ($12$h), S ($13$h), R ($14$h), Q ($16$h). P is kept; S starts at $11$h, before $12$h: no; R starts at
$13$h: kept; Q starts at $12$h, before $14$h: no. P and R, two bookings, as many as possible.
\end{solution}

\begin{solution}{exo:g11:greedy-algorithms:4}
Spices ($2$~kg, $50$), tea ($4$~kg, $60$), then $2$ of the $6$~kg of coffee, worth $72 \times 2/6 = 24$: in all
$134$.
\end{solution}

\begin{solution}{exo:g11:greedy-algorithms:5}
\smallskip\noindent\hspace*{2em}\code{def greedy_change(amount, coins):}

\noindent\hspace*{4em}\code{given = []}

\noindent\hspace*{4em}\code{for c in sorted(coins, reverse=True):}

\noindent\hspace*{6em}\code{k, amount = divmod(amount, c)}

\noindent\hspace*{6em}\code{given += [c] * k}

\noindent\hspace*{4em}\code{return given}

\smallskip\noindent \code{divmod(a, c)} gives the number of coins \code{c} that fit and what is left, in one step.
\end{solution}

\begin{solution}{exo:g11:greedy-algorithms:6}
L from $8$h to $20$h, x from $9$h to $10$h, y from $10$h to $11$h, z from $11$h to $12$h. By start, L comes first and
is kept; it overlaps all the others. Earliest end first keeps x, y and z.
\end{solution}

\begin{solution}{exo:g11:greedy-algorithms:7}
\smallskip\noindent\hspace*{2em}\code{def max_activities(requests):}

\noindent\hspace*{4em}\code{chosen, free_from = [], 0}

\noindent\hspace*{4em}\code{for name, start, end in sorted(requests, key=lambda r: r[2]):}

\noindent\hspace*{6em}\code{if start >= free_from:}

\noindent\hspace*{8em}\code{chosen.append(name)}

\noindent\hspace*{8em}\code{free_from = end}

\noindent\hspace*{4em}\code{return chosen}

\smallskip\noindent The tests compare its count with the best found by trying every subset, on $200$ random sets of
requests. (The \code{0} start assumes hours are not negative.)
\end{solution}

\begin{solution}{exo:g11:greedy-algorithms:8}
With $k$ coins of $5$, from $0$ to $4$, the rest $20 - 5k$ is paid with $1$ and $2$, trying every number of coins of
$2$: $11$, $8$, $6$, $3$ and $1$ ways for $k = 0$ to $4$. In all $29$, as printed.
\end{solution}

\begin{solution}{exo:g11:greedy-algorithms:9}
By value per kilogram: spices, tea, and $4$~kg of coffee, $50 + 60 + 48 = 158$. By value: coffee ($72$, $6$~kg) and
tea ($60$, $4$~kg), $132$. Per kilogram is better: the sack's limit is in kilograms, and the spices, worth less in all,
give the most for each kilogram they use.
\end{solution}

\begin{solution}{exo:g11:greedy-algorithms:10}
Values per kilogram: X $8$, Y and Z $7$. Greedy keeps X ($5$~kg), and then neither Y nor Z fits: value $40$. Y and Z
together weigh $8$~kg and are worth $56$, which is optimal.
\end{solution}

\begin{solution}{exo:g11:greedy-algorithms:11}
\smallskip\noindent\hspace*{2em}\code{def is_canonical_upto(coins, m):}

\noindent\hspace*{4em}\code{return all(len(greedy_change(a, coins)) == fewest(a, sorted(coins))[0]}

\noindent\hspace*{10em}\code{for a in range(1, m + 1))}

\smallskip\noindent By the theorem, $m = c_m + c_{m-1} - 1$, the largest coin plus the next one, minus one, is
enough: if no amount up to it fails, none fails at all.
\end{solution}

\begin{solution}{exo:g11:greedy-algorithms:12}
Let the requests be sorted by end hour, and let $o_1, o_2, \dots, o_p$ be an optimal choice, in order of time. Invariant of
the loop: after greedy has kept $g_1, \dots, g_k$, with $k \le p$, each $g_i$ ends no later than $o_i$. It holds for
$k = 0$. When greedy has $g_1, \dots, g_k$ and $k < p$, the request $o_{k+1}$ starts after $o_k$ ends, hence after
$g_k$ ends: it is compatible with what greedy kept, and greedy, scanning by end hour, keeps a request no later than
$o_{k+1}$ in that order, so its $g_{k+1}$ ends no later than $o_{k+1}$. The invariant is kept as long as $k < p$, so
greedy keeps at least $p$ requests, and therefore exactly $p$, the most possible. The loop is a \code{for} over the
sorted list and stops.
\end{solution}

\begin{solution}{exo:g11:greedy-algorithms:13}
Take a way of paying an amount with the fewest coins: $f$ coins of $5$, $t$ of $2$, $u$ of $1$. Then $u \le 1$, since
$1 + 1$ could become $2$; $t \le 2$, since $2 + 2 + 2$ could become $5 + 1$; and not $t = 2$ with $u = 1$, since
$2 + 2 + 1$ could become $5$. So $2t + u$ is one of $0$, $1$, $2$, $2 + 1 = 3$, $2 + 2 = 4$: it is less than $5$, $f$ is
the number of fives that fit, and $2t + u$ is the remainder. Each remainder from $0$ to $4$ has exactly one such way,
and it is greedy's: greedy also takes every five that fits, then a $2$ while it fits, then a $1$.
\end{solution}

\begin{solution}{exo:g11:greedy-algorithms:14}
\smallskip\noindent\hspace*{2em}\code{def rooms_needed(requests):}

\noindent\hspace*{4em}\code{rooms = []}

\noindent\hspace*{4em}\code{for _, start, end in sorted(requests, key=lambda r: r[1]):}

\noindent\hspace*{6em}\code{for i, free in enumerate(rooms):}

\noindent\hspace*{8em}\code{if free <= start:}

\noindent\hspace*{10em}\code{rooms[i] = end}

\noindent\hspace*{10em}\code{break}

\noindent\hspace*{6em}\code{else:}

\noindent\hspace*{8em}\code{rooms.append(end)}

\noindent\hspace*{4em}\code{return len(rooms)}

\smallskip\noindent \code{rooms[i]} is the hour at which room \code{i} becomes free; the \code{else} of a \code{for}
runs when the loop ended without \code{break}. A new room is opened only when every room is busy at that start: at that
hour, as many requests as rooms, plus the new one, are going on, so no plan can use fewer rooms. The tests check the
equality on $200$ random sets.
\end{solution}

\begin{solution}{exo:g11:greedy-algorithms:15}
P from $8$h to $12$h, S from $11$h to $13$h, Q from $12$h to $16$h. Shortest first keeps S ($2$~h), which overlaps P and
Q: one booking. P and Q together make two.
\end{solution}

\begin{solution}{pb:g11:greedy-algorithms:1}
\textbf{1.} $24 = 12 + 12$; $30 = 25 + 5$; $37 = 25 + 12$: two coins each.

\textbf{2.} Without a coin of $1$, some amounts cannot be paid at all, and greedy may leave a remainder smaller than
every coin, as $3$ with coins of $2$ and $5$ after taking nothing, or $1$ after taking a $2$.

\textbf{3.} Below $c_m + c_{m-1} = 40 + 25 = 65$.

\textbf{4.} The theorem places the first failure; the question asks for all of them up to $100$.

\textbf{5.} Greedy: $25 + 5 + 5 + 1$, four coins. Better: $12 + 12 + 12$, three.

\textbf{6.} The first failure is above $c_3 + 1 = 6$: amounts $1$ to $6$ cannot be it. The program examined $7$ to $36$
(and, without knowing the theorem, $1$ to $6$ too) before finding $36$.

\textbf{7.} $49$: greedy $40 + 5 + 2 + 2$, four coins; $25 + 12 + 12$, three. $50$: greedy $40 + 5 + 5$, three;
$25 + 25$, two.

\textbf{8.} One division per coin value, six in all, whatever the amount; the search tries combinations of all the
coins.

\textbf{9.} No: $36$ fails.

\textbf{10.} Their best change uses three coins of $12$: $36 = 3 \times 12$, $61 = 25 + 3 \times 12$,
$76 = 40 + 3 \times 12$. Greedy takes a larger coin first and can no longer use three $12$.

\textbf{11.} $60$: greedy $40 + 12 + 5 + 2 + 1$, five; $25 + 25 + 5 + 5$, four. $61$: greedy $40 + 12 + 5 + 2 + 2$,
five; $25 + 12 + 12 + 12$, four. $62$: greedy $40 + 12 + 5 + 5$, four; $25 + 25 + 12$, three. $63$: greedy
$40 + 12 + 5 + 5 + 1$, five; $25 + 25 + 12 + 1$, four.

\textbf{12.} Greedy gives $40 + 40 + 12 + 5 + 2 + 1$, six coins; $25 + 25 + 25 + 25$ needs four: two too many. $100$
is the only amount up to $100$ where greedy loses two coins.

\textbf{13.} $17$.

\textbf{14.} Up to $40 + 25 = 65$; by the theorem, a system that does not fail below that never fails.

\textbf{15.} Greedy gives $25 + 10 + 1$, three coins; two coins cannot make $36$ ($40$ is too large, and no two of
$25$, $10$, $5$, $2$, $1$ add to $36$).

\textbf{16.} $36$: greedy $26 + 5 + 5$, three, as many as $12 \times 3$. $50$: greedy $40 + 5 + 5$, three; with $25$
gone, no two coins make $50$.

\textbf{17.} A coin already present would leave five different coins: the system is not one coin changed, but one coin
removed.

\textbf{18.} Replacing $25$ by $3$, by $19$ or by $26$.

\textbf{19.} Shopkeepers can follow the quickest rule and still give the fewest coins, so fewer coins change hands and
fewer need to be minted.

\textbf{20.} In Brenmoor greedy change \textbf{first goes wrong at $36$}, giving four coins where three of $12$ suffice,
and it goes wrong on \textbf{$17$ amounts up to $100$}.
\end{solution}
